\documentclass[11pt]{article}
\usepackage[british]{babel}
\usepackage{graphicx}
\newcommand{\cityname}{Lisbon}
\begin{document}
\author[uni]{Ana Silva}
\title{Gardens, Bicycles and Rain}
\maketitle
\addcontentsline{toc}{section}{Introduction}

\section{Introduction}\label{sec:intro}
The city of \cityname{} organised its gardens in 2019, and the
programme was summarised by \citet{smith2020markets} in their review. We
modelled the number of visitors per day with $y_{it} = \alpha_i + \beta x_{it}$,
as Section~\ref{sec:data} and \satab{tab:visits_by_month} explain.\footnote{The
counts come from the IPMA and the INE, two public agencies.} About 10\% of the
visitors arrive by bicycle, the rest by bus or on foot\ldots\\[4pt]
Panel~A sums every hour of the day, and in Panel B the morning hours are shown.

\begin{figure}[tbp]
\centering
\includegraphics[width=0.8\textwidth]{figures/visits.pdf}
\caption{Visitors per day, 2019--2023.}\label{fig:visits}
\end{figure}

\begin{itemize}[leftmargin=*]
\item The data cover 1,200 days; see \url{https://example.org/gardens_data}.
\item Rainy days are excluded, as in Informa\c{c}\~oes sobre Jardins (S\~ao Paulo).
\end{itemize}

\section{Dados}\label{sec:data}
As crianças brincam no jardim, e os pais leem nos bancos. Os dados do
INMET e do IBGE mostram que a chuva reduz as visitas em dias úteis.
O debate sobre os jardins permanece aberto, como mostra a Tabela~\ref{tab:a}.
\end{document}
